\documentclass[../../main.tex]{subfiles}
\begin{document}
\begin{oframed}
  \subfile{problem}
\end{oframed}

{\bf [解答]}

題意より
\begin{align}
  a > 0 \label{eq:1}
\end{align}
とする．題意の二つの放物線を
\begin{align}
  \begin{cases}
    C_1 : y = x^2 \\
    C_2 : x = \dfrac{1}{6}a^3 y^2 - \dfrac{7}{6}ay \equiv f(y)
  \end{cases}
\end{align}
とおく．$C_2$を平方完成すると
\begin{align}
  f(y) = \dfrac{1}{6}a^3 \left( y - \dfrac{7}{2a^2} \right)^2 - \dfrac{49}{24a}
\end{align}
である．

まずは二つの放物線の交点の$x$座標を求める．$C_1, C_2$から$y$を消去して
\begin{align}
  & 6x = a^3 x^4 - 7ax^2 \\
  & x(a^3 x^3 - 7ax - 6) = 0 \\
  & x(ax + 1)((ax)^2 - (ax) - 6) = 0 \\
  & x(ax + 1)(ax - 3)(ax + 2) = 0 \\
  \therefore
  & x = 0, \, -\dfrac{1}{a}, \, -\dfrac{2}{a}, \, \dfrac{3}{a}
\end{align}
となる．
よってグラフの概形は下図のようになる．

\begin{figure}[htb]
  \centering
  \begin{tikzpicture}[scale=1.2, >=stealth]
    \fill[pattern=north east lines, pattern color=blue!75!black]
    (0,0) -- plot[domain=0:1, samples=40, variable=\y]
    ({(\y*\y-7*\y)/6},\y)
    -- plot[domain=1:0, samples=40, variable=\y] ({-sqrt(\y)},\y) -- cycle;
    \fill[pattern=north west lines, pattern color=teal!80!black]
    (-1,1) -- plot[domain=1:4, samples=60, variable=\y]
    ({(\y*\y-7*\y)/6},\y)
    -- plot[domain=4:1, samples=60, variable=\y] ({-sqrt(\y)},\y) -- cycle;
    \draw[->] (-2.5, 0) -- (2.0, 0) node[right] {$x$};
    \draw[->] (0, -0.5) -- (0, 4.5) node[above] {$y$};
    \node[below left] at (0,0) {$0$};

    \draw[thick, domain=-2.1:1.8, samples=50] plot (\x, {\x*\x}) node[above] {$C_1$};

    \draw[thick, domain=0:4.2, samples=50] plot ({1/6*\x*\x - 7/6*\x}, \x) node[right] {$C_2$};

    \draw[dashed] (-2.5, 4) -- (0, 4) node[right] {$\frac{4}{a^2}$};
    \draw[densely dotted] (-49/24, 3.5) -- (0, 3.5) node[right] {$\frac{7}{2a^2}$};
    \draw[dashed] (-2.5, 1) -- (0, 1) node[right] {$\frac{1}{a^2}$};
    \draw[densely dotted] (-49/24, 0) node[above left] {$-\frac{49}{24a}$} -- (-49/24, 3.5);

    \draw[dashed] (-2, 0) node[below] {$-\frac{2}{a}$} -- (-2, 4);
    \draw[dashed] (-1, 0) node[below] {$-\frac{1}{a}$} -- (-1, 1);

    \node[circle, fill=black, inner sep=1.3pt] at (0,0) {};
    \node[circle, fill=black, inner sep=1.3pt] at (-1,1) {};
    \node[circle, fill=black, inner sep=1.3pt] at (-2,4) {};
    \node[circle, fill=black, inner sep=1.3pt] at (-49/24,3.5) {};
    \node at (-2.35,2.45) {$S_1$};
    \draw[->] (-2.13,2.38) -- (-1.72,2.38);
    \node at (0.35,0.62) {$S_2$};
    \draw[->] (0.17,0.55) -- (-0.66,0.55);
  \end{tikzpicture}
  \caption{二つの放物線$C_1,C_2$とその作る面積$S_1,S_2$．}
\end{figure}

\begin{figure}[htb]
  \centering
  \begin{tikzpicture}[x=1.4cm, y=0.95cm, >=stealth, every node/.style={font=\small}]
    \node at (-1,3.65) {$S_2\ (0\le y\le1)$};
    \node at (4,3.65) {$S_1\ (1\le y\le4)$};
    \begin{scope}[xscale=2, yscale=3]
      % 左：長方形を C_1 と C_2 で三つに分ける
      \fill[pattern=north east lines, pattern color=gray!70!black]
      (-1,0) -- (0,0) -- plot[domain=0:-1, samples=50] (\x,{\x*\x}) -- cycle;
      \fill[pattern=north west lines, pattern color=blue!75!black]
      (0,0) -- plot[domain=0:1, samples=50, variable=\y]
      ({(\y*\y-7*\y)/6},\y)
      -- plot[domain=1:0, samples=50, variable=\y] ({-sqrt(\y)},\y) -- cycle;
      \fill[pattern=north east lines, pattern color=gray!70!black]
      (0,0) -- plot[domain=0:1, samples=50, variable=\y]
      ({(\y*\y-7*\y)/6},\y) -- (0,1) -- cycle;
      \draw[thick] (-1,0) rectangle (0,1);
      \draw[very thick, blue!85!black] plot[domain=-1:0, samples=60] (\x,{\x*\x});
      \draw[very thick, teal!85!black] plot[domain=0:1, samples=60, variable=\y]
      ({(\y*\y-7*\y)/6},\y);
    \end{scope}
    \draw[->, thick] (-2.25,0) -- (0.25,0) node[right] {$x$};
    \draw[->, thick] (0,-0.2) -- (0,3.25) node[above] {$y$};
    \node[left] at (-2,0) {$-1$};
    \node[right] at (0,3) {$1$};
    \node[blue!85!black] at (-2.1,3.25) {$C_1$};
    \node[teal!85!black] at (-0.25,3.25) {$C_2$};
    \node at (0.6,1.6) {$S_2$};
    \draw[->] (0.35,1.5) -- (-1.32,1.5);
    \begin{scope}[shift={(5,-1)}]
      % 右：C_2 と y 軸の間から C_1 と y 軸の間を差し引く
      \fill[pattern=north east lines, pattern color=teal!80!black]
      (0,1) -- plot[domain=1:4, samples=70, variable=\y]
      ({(\y*\y-7*\y)/6},\y) -- (0,4) -- cycle;
      \fill[pattern=north east lines, pattern color=gray!70!black]
      (0,1) -- (-1,1) -- plot[domain=1:4, samples=70, variable=\y]
      ({-sqrt(\y)},\y) -- (0,4) -- cycle;
      \draw[very thick, blue!85!black] plot[domain=1:4, samples=70, variable=\y]
      ({-sqrt(\y)},\y);
      \draw[very thick, teal!85!black] plot[domain=1:4, samples=70, variable=\y]
      ({(\y*\y-7*\y)/6},\y);
    \end{scope}
    \draw[densely dotted, thick] (3,0) -- (5,0);
    \draw[densely dotted, thick] (3,3) -- (5,3);
    \draw[->, thick] (5,-0.2) -- (5,3.25) node[above] {$y$};
    \node[right] at (5,0) {$1$};
    \node[right] at (5,3) {$4$};
    \node[blue!85!black] at (4.6,3.25) {$C_1$};
    \node[teal!85!black] at (2.9,3.25) {$C_2$};
    \node at (2.5,1.65) {$S_1$};
    \draw[->] (2.7,1.55) -- (3.25,1.55);
  \end{tikzpicture}
  \caption{$a=1$ とした放物線上での面積の構成（青線が $C_1$，緑線が $C_2$）．左は長方形から灰色の二領域を除くと $S_2$，右は $C_2$ と $y$ 軸の間から灰色の領域を除くと $S_1$．}
\end{figure}

二つの放物線の囲む二つの部分の面積を図のように$S_1, S_2$とおく．
$S_2$は二辺の長さが$\dfrac{1}{a^2}, \dfrac{1}{a}$の長方形の面積から$C_1, C_2$の作る二つの部分の面積を減じたものである．

\begin{align}
  S_2 =
  \begin{tikzpicture}[scale=0.4, baseline={([yshift=-0.5ex]current bounding box.center)}]
    \fill[pattern=north east lines, pattern color=gray!70!black] (-1,0) rectangle (0,1);
    \draw (-1,0) rectangle (0,1);
  \end{tikzpicture}
  -
  \begin{tikzpicture}[scale=0.4, baseline={([yshift=-0.5ex]current bounding box.center)}]
    \fill[pattern=north east lines, pattern color=blue!75!black]
    (-1,0) -- (0,0) -- plot[domain=0:-1, samples=20] (\x,{\x*\x}) -- cycle;
    \draw (-1,0) -- (0,0) -- plot[domain=0:-1, samples=20] (\x,{\x*\x}) -- cycle;
  \end{tikzpicture}
  -
  \begin{tikzpicture}[scale=0.4, baseline={([yshift=-0.5ex]current bounding box.center)}]
    \fill[pattern=north east lines, pattern color=teal!80!black]
    (0,0) -- plot[domain=0:1, samples=20, variable=\y]
    ({(\y*\y-7*\y)/6},\y) -- (0,1) -- cycle;
    \draw (0,0) -- plot[domain=0:1, samples=20, variable=\y]
    ({(\y*\y-7*\y)/6},\y) -- (0,1) -- cycle;
  \end{tikzpicture}
  \label{eq:2}
\end{align}
となる．各項を計算すると
\begin{align}
  \begin{tikzpicture}[scale=0.4, baseline={([yshift=-0.5ex]current bounding box.center)}]
    \fill[pattern=north east lines, pattern color=gray!70!black] (-1,0) rectangle (0,1);
    \draw (-1,0) rectangle (0,1);
  \end{tikzpicture}
  &= \dfrac{1}{a} \cdot \dfrac{1}{a^2} = \dfrac{1}{a^3} \\
  \begin{tikzpicture}[scale=0.4, baseline={([yshift=-0.5ex]current bounding box.center)}]
    \fill[pattern=north east lines, pattern color=blue!75!black]
    (-1,0) -- (0,0) -- plot[domain=0:-1, samples=20] (\x,{\x*\x}) -- cycle;
    \draw (-1,0) -- (0,0) -- plot[domain=0:-1, samples=20] (\x,{\x*\x}) -- cycle;
  \end{tikzpicture}
  &= \int_0^{\frac{1}{a}} x^2 \, dx = \dfrac{1}{3}\dfrac{1}{a^3} \\
  \begin{tikzpicture}[scale=0.4, baseline={([yshift=-0.5ex]current bounding box.center)}]
    \fill[pattern=north east lines, pattern color=teal!80!black]
    (0,0) -- plot[domain=0:1, samples=20, variable=\y]
    ({(\y*\y-7*\y)/6},\y) -- (0,1) -- cycle;
    \draw (0,0) -- plot[domain=0:1, samples=20, variable=\y]
    ({(\y*\y-7*\y)/6},\y) -- (0,1) -- cycle;
  \end{tikzpicture}
  &= -\int_0^{\frac{1}{a^2}} f(y) \, dy = -\left[ \dfrac{a^3}{18}y^3 - \dfrac{7}{12}ay^2 \right]_0^{\frac{1}{a^2}} = \dfrac{19}{36}\dfrac{1}{a^3}
\end{align}
だから，\cref{eq:2}に代入して
\begin{align}
  S_2 = \dfrac{\left( 1 - \dfrac{1}{3} - \dfrac{19}{36} \right)}{a^3} = \dfrac{5}{36a^3} \label{eq:3}
\end{align}
を得る．

一方 $S_1$ は$\dfrac{1}{a^2}\le y \le \dfrac{4}{a^2}$での$C_1, C_2$の作る領域の差分だから
\begin{align}
  S_1 =
  \begin{tikzpicture}[scale=0.4, baseline={([yshift=-0.5ex]current bounding box.center)}]
    \fill[pattern=north east lines, pattern color=teal!80!black]
    (0,1) -- plot[domain=1:4, samples=30, variable=\y]
    ({(\y*\y-7*\y)/6},\y) -- (0,4) -- cycle;
    \draw (0,1) -- plot[domain=1:4, samples=30, variable=\y]
    ({(\y*\y-7*\y)/6},\y) -- (0,4) -- cycle;
  \end{tikzpicture}
  -
  \begin{tikzpicture}[scale=0.4, baseline={([yshift=-0.5ex]current bounding box.center)}]
    \fill[pattern=north east lines, pattern color=blue!75!black]
    (0,1) -- (-1,1) -- plot[domain=1:4, samples=30, variable=\y]
    ({-sqrt(\y)},\y) -- (0,4) -- cycle;
    \draw (0,1) -- (-1,1) -- plot[domain=1:4, samples=30, variable=\y]
    ({-sqrt(\y)},\y) -- (0,4) -- cycle;
  \end{tikzpicture}
  \label{eq:4}
\end{align}
である．$y=x^2$の逆関数が$x=\pm \sqrt{y}$であることに注意して各項計算すると
\begin{align}
  \begin{tikzpicture}[scale=0.4, baseline={([yshift=-0.5ex]current bounding box.center)}]
    \fill[pattern=north east lines, pattern color=teal!80!black]
    (0,1) -- plot[domain=1:4, samples=30, variable=\y]
    ({(\y*\y-7*\y)/6},\y) -- (0,4) -- cycle;
    \draw (0,1) -- plot[domain=1:4, samples=30, variable=\y]
    ({(\y*\y-7*\y)/6},\y) -- (0,4) -- cycle;
  \end{tikzpicture}
  &= -\int_{\frac{1}{a^2}}^{\frac{4}{a^2}} f(y) \, dy = -\left[ \dfrac{a^3}{18}y^3 - \dfrac{7}{12}ay^2 \right]_{\frac{1}{a^2}}^{\frac{4}{a^2}} = \dfrac{21}{4a^3} \\
  \begin{tikzpicture}[scale=0.4, baseline={([yshift=-0.5ex]current bounding box.center)}]
    \fill[pattern=north east lines, pattern color=blue!75!black]
    (0,1) -- (-1,1) -- plot[domain=1:4, samples=30, variable=\y]
    ({-sqrt(\y)},\y) -- (0,4) -- cycle;
    \draw (0,1) -- (-1,1) -- plot[domain=1:4, samples=30, variable=\y]
    ({-sqrt(\y)},\y) -- (0,4) -- cycle;
  \end{tikzpicture}
  &= \int_{\frac{1}{a^2}}^{\frac{4}{a^2}} \sqrt{y} \, dy = \left[ \dfrac{2}{3} y^{\frac{3}{2}}\right]_{\frac{1}{a^2}}^{\frac{4}{a^2}} = \dfrac{14}{3a^3}
\end{align}
だから，\cref{eq:4}に代入して
\begin{align}
  S_1 = \dfrac{\left( \dfrac{21}{4} - \dfrac{14}{3} \right)}{a^3} = \dfrac{7}{12a^3} \label{eq:5}
\end{align}
を得る．
\cref{eq:3,eq:5}から求める面積比は
\begin{align}
  S_1 : S_2 = \dfrac{7}{12} : \dfrac{5}{36} = 21 : 5
\end{align}
である．

\newpage

\end{document}
